\documentclass[russian]{beamer}
\usepackage[T2A]{fontenc}
\usepackage[english,main=russian]{babel}
\usepackage{
    amssymb,mathpartir,mathtools,simplebnf,stmaryrd,
    tikz-cd
}
\title{Чистая система типов с рангами функций в качестве сортов}
\author{Соколов П.П.}
\institute{Высшая Школа Экономики}
\date{2026}
\begin{document}
\frame{\titlepage}
\begin{frame}{План на сегодня}
\tableofcontents
\end{frame}
\section{Мотивация и идея}
\begin{frame}{Мотивация и идея}
\begin{itemize}
\item Больше всего вопросов при изучении современных
    средств интерактивных доказательств (Lean, Rocq,
    Agda...) вызывает устройство \textit{иерархий
    вселенных}: $\mathtt{Type}_1 : \mathtt{Type}_2 :
    \ldots$
\item Вызывают они вопросы и с метатеоретической точки
    зрения --- в доказательствах здравости системы
    $\mathtt{Type}_n$ интерпретируется как $n$-й
    недостижимый кардинал.
\item Поэтому интересно искать альтернативные варианты
    устройства иерархии. Далее описан забавный вариант,
    располагающий типы по ``вселенным'' соответственно
    \textit{рангу функций}.
\item Б\'{о}льшая часть описанных результатов формально
    верифицирована на Rocq $\checkmark$.
\end{itemize}
\end{frame}
\section{Чистые системы типов}
\begin{frame}{Чистые системы типов: синтаксис
    $\checkmark$}
\begin{bnf}
$x$ : \textsf{Переменная} :in: $X$;;
$s$ : \textsf{Сорт} :in: $S$;;
$T$, $U$, $t$, $u$ : \textsf{Терм} ::=
| $x$ : переменная
| $s$ : сорт как тип
| $\Pi(x\colon T).\,U$ : тип зависимых ф-й
| $\lambda(x\colon T).\,u$ : анонимная функция
| $u\,t$ : применение функции
;;
\end{bnf}
\end{frame}
\begin{frame}{ЧСТ: правила одношаговой редукции
    $\checkmark$}
\begin{mathpar}
\inferrule{T \to_\beta T'}
    {\Pi(x:T).U \to_\beta \Pi(x:T').U}
\and
\inferrule{U \to_\beta U'}
    {\Pi(x:T).U \to_\beta \Pi(x:T).U'}
\and
\inferrule{T \to_\beta T'}
    {\lambda(x:T).t \to_\beta \lambda(x:T').t}
\and
\inferrule{t \to_\beta t'}
    {\lambda(x:T).t \to_\beta \lambda(x:T).t'}
\and
\inferrule{t \to_\beta t'}{t\,u \to_\beta t'\,u}
\and
\inferrule{u \to_\beta u'}{t\,u \to_\beta t\,u'}
\and
\inferrule{}
    {(\lambda(x:T).t)\,u \to_\beta t[x\coloneqq u]}
\end{mathpar}
\end{frame}
\begin{frame}{ЧСТ: правила типизации $\checkmark$}
\begin{mathpar}
\inferrule{x:T\in\Gamma}{\Gamma\vdash x:T}
\and
\inferrule{ }{\Gamma\vdash s:s'} (s, s')\in\mathrm{Ax}
\and
\inferrule{
    \Gamma\vdash T:s_T \\
    \Gamma,x:T\vdash U:s_U
}{\Gamma\vdash\Pi(x:T).\,U:s}
(s_T, s_U, s)\in\mathrm{Rule}
\and
\inferrule{
    \Gamma\vdash T:s_T \\
    \Gamma,x:T\vdash u:U
}{\Gamma\vdash\lambda(x:T).u:\Pi(x:T).U}
\and
\inferrule{
    \Gamma\vdash u:\Pi(x:T).U \\
    \Gamma\vdash t:T
}{\Gamma\vdash u\,t:U[x\coloneqq t]}
\and
\inferrule{
    \Gamma\vdash t:T \\
    \Gamma\vdash U:s \\
    T =_\beta U
}{\Gamma\vdash t:U}
\end{mathpar}
\end{frame}
\section{<<Комбинаторные>> свойства ЧСТ}
\begin{frame}[fragile]
\frametitle{Метатеоретические свойства ЧСТ}
\begin{itemize}
    \item $\checkmark$ Отношение многошаговой редукции
        $(\to_\beta^*)$ конфлюэнтно, т.е.
        \[\begin{tikzcd}
            t \arrow[r] \arrow[d]
            & t' \arrow[d, dashed] \\
            t'' \arrow[r, dashed] & t^*
        \end{tikzcd}\]
    \item $\checkmark$ Редукция сохраняет типизацию,
        т.е.
        \[\dfrac{\Gamma\vdash t:T\quad t\to_\beta t'}
                {\Gamma\vdash t':T}\]
    \item Здравы ли ЧСТ? (<<есть ли полезная модель у
        ЧСТ?>> или <<все ли типы населены?>> или
        <<любой ли терм нормализуется?>>)
\end{itemize}
\end{frame}
\section{О здравости ЧСТ}
\begin{frame}{Какие ЧСТ здравы?}
\begin{itemize}
    \item Исчисление конструкций:
        \[\begin{array}{c}
            S=\{*,\square\}
            \quad \mathrm{Ax}=\{(*,\square)\}
            \\ \mathrm{Rule}=\{
                (*,*,*),(*,\square,\square),
                (\square,*,*),(\square,\square,\square)
            \}
        \end{array}\]
    \item Первая счётная иерархия \textit{вселенных}:
        \[\begin{array}{c}
            S=\{\mathcal{U}_\ell\;|\;\ell\in\mathbb{N}\}
            \quad \mathrm{Ax}=\{
                (\mathcal{U}_\ell,\mathcal{U}_{\ell+1})
                \;|\;\ell\in\mathbb{N}
            \}
            \\ \mathrm{Rule}=\{
                (\mathcal{U}_\ell, \mathcal{U}_{\ell'},
                \mathcal{U}_{\max(\ell,\ell')}) \;|\;
                \ell,\ell'\in\mathbb{N}
            \}
        \end{array}\]
    \item Счётная иерархия с импредикативным
        \texttt{Prop}:
        \[\begin{array}{rcl}
            S &=&\{\mathtt{Prop}\}\sqcup
            \{\mathcal{U}_\ell\;|\;\ell\in\mathbb{N}\}
            \\ \mathrm{Ax} &=&
            \{(\mathtt{Prop},\mathcal{U}_0)
            \}\sqcup\{
                (\mathcal{U}_\ell,\mathcal{U}_{\ell+1})
                \;|\;\ell\in\mathbb{N}
            \}
            \\ \mathrm{Rule} &=&\{
                (\mathcal{U}_\ell,\mathtt{Prop},
                \mathtt{Prop}) \;|\;
                \ell\in\mathbb{N}
            \} \;\sqcup
            \\ & \sqcup & \{
                (\mathcal{U}_\ell, \mathcal{U}_{\ell'},
                \mathcal{U}_{\max(\ell,\ell')}) \;|\;
                \ell,\ell'\in\mathbb{N}
            \}
        \end{array}\]
\end{itemize}
\end{frame}
\begin{frame}{Парадокс Жирара}
    \begin{theorem}[Girard'72]
    Следующая ЧСТ (известная как System U) содержит
    \textbf{не нормализующийся} терм:
    \[\begin{array}{c}
        S=\{*,\square,\triangle\}
        \quad\mathrm{Ax}=\{
            (*,\square),(\square,\triangle)
        \}
        \\ \mathrm{Rule}=\{
            (*,*,*),(\square,*,*),(\triangle,*,*),
            (\square,\square,\square),
            (\triangle,\square,\square)
        \}
    \end{array}\]
    Более того, этот терм имеет тип $\Pi(T:*).T$.
    \end{theorem}
    Другими словами, <<здравая ЧСТ не может поддерживать
    одновременно \textbf{и} импредикативный
    \texttt{Prop}, \textbf{и} полиморфные функции>>.

    Из парадокса Жирара выводится противоречивость ещё
    более простой ЧСТ, известной как Type-in-Type:
    \[S=\{*\}\quad\mathrm{Ax}=\{(*,*)\}\quad
    \mathrm{Rule}=\{(*,*,*)\}\]
\end{frame}
\begin{frame}{Как насчёт рангов?}
\textit{Рангом} функции называется максимальная
вложенность $\Pi$ в позиции аргумента:
\[\mathrm{rk}\,(\Pi(x:T).U)=\max\,(\mathrm{rk}\,T+1,
    \mathrm{rk}\,U)\]
Рассмотрим $S=\{\mathrm{rk}_n\;|\;n\in\mathbb{N}\}$, и
\[\mathrm{Rule}=\{(\mathrm{rk}_m,\mathrm{rk}_n,
    \mathrm{rk}_{\max(m+1,n)})\;|\;m,n\in\mathbb{N}\}.\]
В качестве $\mathrm{Ax}$ можно взять цепочку
$\mathrm{rk}_0:\mathrm{rk}_1:\mathrm{rk}_2:\ldots$, но
получится система, строго менее удобная, чем обычная
иерархия вселенных. Попробуем поступить следующим
образом:
\[\mathrm{Ax}=\{(\mathrm{rk}_n,\mathrm{rk}_0)\;|\;
    n\in\mathbb{N}\}\]
В частности, $\mathrm{rk}_0:\mathrm{rk}_0$. Но в данной
системе парадокс Жирара не применим!
\end{frame}
\begin{frame}{Варианты доказательств здравости ЧСТ}
\begin{itemize}
    \item Наличие хорошей модели;
    \item Normalization-by-evaluation: наличие модели,
        для которой есть операция <<квотинга>>, по
        значению в модели получающая соответствующую
        нормальную форму;
    \item Доказательство строгой нормализации;
    \item \dots
\end{itemize}
\end{frame}
\begin{frame}{<<Общее доказательство нормализации для
    ЧСТ>>}
\begin{itemize}
    \item $\checkmark$ (насыщенное) $\Lambda$-множество:
        множество \textit{реализаторов} $S$ вместе с
        отношением
        $(\vDash)\subseteq S\times\mathrm{Term}$ (и
        дополнительные условия, связанные с реализацией
        строго нормализуемых термов);
    \item $\checkmark$ $\mathfrak{E}$-множество:
        семейство $\Lambda$-множеств $F$ вместе с
        отношениями эквивалентности $\|i\|$ на $F$ и
        $|i|$ на $\cup F$ для каждого
        $i\in\mathfrak{E}$.
    \item $\checkmark$ Иерархия вселенных:
        изоморфные друг другу $\mathfrak{E}$-множество
        $\mathfrak{A}^\Uparrow(s)$ и $\Lambda$-множество
        $\mathfrak{A}^\Downarrow(s)$ для каждого сорта
        $s$ с рядом допусловий.
    \item Пара интерпретаций
        $[\Gamma\vdash t](\gamma)$ и $\llbracket\Gamma
        \vdash t\rrbracket(\gamma)$ такие, что, если
        $\Gamma\vdash T:s$, то
        $[\Gamma\vdash T]\in\mathfrak{A}^\Downarrow(s)$
        и $\llbracket\Gamma\vdash T\rrbracket\in
        \mathfrak{A}^\Uparrow(s)$. Из корректности
        интерпретации следует сильная нормализация.
    \item Собственно, конструкция иерархии для нашей
        ЧСТ.
\end{itemize}
\end{frame}
\end{document}
